\scanpage{3}
$k$ corps de car $p\unsure{\mathord{-}|>}0$

$G$ gpe $p$-divisible, $G^{*}$ dual

$g=\dim G,\quad g^{*}=\dim G^{*}$

$\widetilde G$ \unsure{gpe|groupe} enveloppant
\[
0\longrightarrow\underline{\mathbb V}(t_{G^{*}})\longrightarrow\widetilde G\longrightarrow G\longrightarrow0
\]

\diagram{\scriptsize\begin{tikzcd}[ampersand replacement=\&,column sep=.38em,row sep=.52em,cells={nodes={inner sep=1pt}}]
\&\&\&\&\&\text{duales}\arrow[ddd,dash]\&\&\&\&\&\\
0\arrow[r]\&t_G^\vee\arrow[r]\&H(G)\arrow[r]\arrow[dd,shift right,"{H(V)=V_0}"']\&t_{G^*}\arrow[r]\&0\&\&0\arrow[r]\&t_{G^*}^\vee\arrow[r]\&t_{\widetilde G}\arrow[r]\arrow[dd,shift right,"{V_0}"']\&t_G\arrow[r]\&0\\
\omega_G\arrow[ur,equal]\&\&\&\&\&\&\omega_{G^*}\arrow[ur,equal]\&\&\&\&\\
0\arrow[r]\&t_G^{\vee(p)}\arrow[r]\&H(G)^{(p)}\arrow[uu,shift right,"{H(F)=\unsure{f_0|F_0}}"']\arrow[r]\&t_{G^*}^{(p)}\arrow[r]\&0\&{}\&0\arrow[r]\&t_{G^*}^{\vee(p)}\arrow[r]\&t_{\widetilde G}^{(p)}\arrow[uu,shift right,"{F_0}"']\arrow[r,dash]\&t_G^{(p)}\arrow[r]\&0\\
\omega_G^{(p)}\arrow[ur,equal]\&\&\&\&\&\&\omega_{G^*}^{(p)}\arrow[ur,equal]\&\&\&\&\\
\&\dim g\&\dim g+g^*\&\dim g^*\&\&\&\&\&\&\&
\end{tikzcd}}

[\uline{NB} $H(G)=t_{\deleted{\unsure{G}}\widetilde{G^{*}}}$]

\textcircled{1} suite exacte pour $t_{\widetilde G}$ et $t_{\widetilde{G^{*}}}$ transposées l'une de l'autre pour un accouplement canonique fonctoriel, \deleted{\unsure{(qui)}} par lequel $F_0$ et $V_0$ (si on est en car $p$) s'interchangent.

\textcircled{2} Dans la suite exacte de gauche, on a
\[
\begin{aligned}
\operatorname{Ker} F_0&=\operatorname{Im} V_0=t_G^{\vee(p)}=\omega_G^{(p)}\\
\operatorname{Ker} V_0&=\operatorname{Im} F_0=\text{\unsure{ce que ça peut} (\unsure{invariant intéressant|nouveaux invariants})}
\end{aligned}
\]

\textcircled{3} $F_0$ et $V_0$ opérant sur $H(G)$ définissent le module de Dieudonné \deleted{\unsure{du gpe} $G$} réduit mod $p$.\margin{?}

\textcircled{4} Le composé $t_{G^*}^{(p)}\xrightarrow{F_0}H(G)\unsure{\longrightarrow|\mathord{-}}t_{G^*}$ correspond : la puissance $p$-ième des dérivations [et correspond : Frobenius du $H^2(G,\unsure{G_{\mathrm a}|G\mathrm a|\mathbb G_m})$]

\deleted{\illegible{3} $\unsure{F_0}$}

\uline{Corollaire}
\[
\begin{aligned}
&G\ \text{purement infinitésimal}&&\Longleftrightarrow F_0\ \text{nilpotent sur }H(G)=t_{\widetilde{G^*}}\\
&\qquad\text{i.e. }G^*\ \text{sans tore}\\
&G^*\ \text{purement infinitésimal}&&\Longleftrightarrow V_0\ \text{nilpotent}\\
&\qquad\text{i.e. }G\ \text{sans composante}\\[-2pt]
&\qquad\qquad\text{torique}\\
&G\ \text{et }G^*\ \text{purement infinitésimaux}&&\Longleftrightarrow F_0\ \text{et }V_0\ \text{nilpotents}
\end{aligned}
\]
% archive sheet number: 2

